\documentclass[11pt]{amsart}
\usepackage{amsmath,amssymb,amsthm}
\usepackage[margin=1.2in]{geometry}
\usepackage{booktabs}
\usepackage{adjustbox}
\usepackage{xcolor}
\usepackage[colorlinks=true,linkcolor=blue,citecolor=blue,urlcolor=blue]{hyperref}
\hypersetup{pdfauthor={Bernd Johannes Wuebben},
  pdftitle={Doubly isolated Batyrev mirror pairs and non-smoothable Calabi-Yau threefolds}}

% Last-resort stretch so a paragraph that cannot break cleanly sets loose
% rather than running into the right margin.
\setlength{\emergencystretch}{2em}

\newtheorem{theorem}{Theorem}[section]
\newtheorem{proposition}[theorem]{Proposition}
\newtheorem{lemma}[theorem]{Lemma}
\newtheorem{corollary}[theorem]{Corollary}
\newtheorem{conjecture}[theorem]{Conjecture}
\theoremstyle{definition}
\newtheorem{definition}[theorem]{Definition}
\newtheorem{example}[theorem]{Example}
\newtheorem{question}[theorem]{Question}
\theoremstyle{remark}
\newtheorem{remark}[theorem]{Remark}

\newcommand{\ZZ}{\mathbb{Z}}
\newcommand{\RR}{\mathbb{R}}
\newcommand{\CC}{\mathbb{C}}
\newcommand{\PP}{\mathbb{P}}
\newcommand{\FF}{\mathbb{F}}
\newcommand{\NN}{N}
\newcommand{\MM}{M}
\newcommand{\conv}{\operatorname{conv}}
\newcommand{\Cone}{\operatorname{Cone}}
\newcommand{\Sing}{\operatorname{Sing}}
\newcommand{\vertt}{\operatorname{vert}}

\title[Doubly isolated Batyrev mirror pairs]{Doubly isolated Batyrev mirror
pairs\\ and non-smoothable Calabi--Yau threefolds}
\author{Bernd Johannes Wuebben}
\subjclass[2020]{14J32 (primary); 14J33, 14M25, 14B07, 52B20 (secondary)}
\keywords{Mirror symmetry, Batyrev duality, reflexive polytopes,
Kreuzer--Skarke classification, Calabi--Yau threefolds, isolated
singularities, smoothability}
\date{August 2, 2026}

\begin{document}

\begin{abstract}
For a reflexive $4$-polytope $\Delta$ with unit edges the generic
anticanonical hypersurface $X \subset \PP_\Delta$ is a Calabi--Yau
threefold with isolated singularities; we ask when the Batyrev mirror
$X^\circ \subset \PP_{\Delta^\circ}$ is isolated-singular as well (we then
call the pair \emph{doubly isolated}).  An exact
identity (the lattice length of the dual edge of a two-dimensional face
equals the multiplicity of its transverse cone) shows this happens if and
only if every two-dimensional face of $\Delta$ is transversally smooth; we
call such $\Delta$ \emph{both-sides unit}.  We determine the both-sides unit
polytopes completely, subject to the database-copy completeness and
no-duplication hypothesis stated in \S\ref{sec:census}.  Under that
hypothesis an exhaustive scan with exact integer predicates and exact
re-verification finds that, among the
$473{,}800{,}776$ reflexive $4$-polytopes of
the Kreuzer--Skarke classification there are exactly $590$, and every
two-dimensional face of every one of them is a triangle, a zonotope, or a
reflexive polygon (one interior point).  Consequently the non-smoothable
germs occurring on doubly isolated mirror pairs are exactly two cyclic
quotient points, $\frac13(1,1,1)$ and $\frac15(1,1,3)$, and the
anticanonical cone over $\FF_1$.  The
germs that deform but admit no smoothing (type \textup{(D)} in the
Altmann--Corti--Filip--Petracci trichotomy, the smallest being the cone over
the Minkowski sum of a triangle and a unit segment) never occur: a Calabi--Yau
threefold with isolated singularities carrying one never has a mirror with
isolated singularities.
At the opposite extreme, exactly one reflexive $4$-polytope yields a
mirror pair with both members \emph{smooth}: the self-dual $24$-cell, whose
hypersurface is the self-mirror threefold with Hodge numbers $(20,20)$.
The $\FF_1$-cone occurs exactly once: a unique pair of polytopes, with $22$
and $26$ vertices, whose hypersurfaces form a mirror pair, with Hodge
numbers $(20,26)$ and $(26,20)$ for the MPCP resolutions, in which $X$
is non-smoothable while
every singular point of $X^\circ$ is locally smoothable.  We also record, as
a further obstruction, that the $26$ node-diagonal vectors of $X^\circ$
admit no all-nonzero relation: any global smoothing would have to couple
the nodal directions to the four del Pezzo-cone directions.  Finally, as a
caution, the natural conjecture suggested by the low-vertex data
(that only triangles and zonotopes occur) is true for all $395{,}406{,}329$
polytopes with at most $18$ vertices and false from $19$ vertices on.
\end{abstract}

\maketitle

\section{Introduction}\label{sec:intro}

Let $\Delta \subset \NN_\RR = \RR^4$ be a reflexive polytope all of whose edges
have lattice length one, let $\Delta^\circ \subset \MM_\RR$ be its polar dual,
and let $X \in |{-K_{\PP_\Delta}}|$ be a generic anticanonical hypersurface in
the Gorenstein Fano toric fourfold associated to the face fan of $\Delta$.
(Throughout, $\Delta$ is the \emph{fan} polytope, Batyrev's $\Delta^\ast$,
so our $\Delta^\circ$ is his $\Delta$; the same convention is in force in
\cite{paper1,paper2}.)  By
Batyrev \cite{Batyrev94}, $X$ is a Calabi--Yau threefold with Gorenstein
canonical singularities, and by
\cite[Prop.~3.1]{paper1} its singular locus is exactly: for each
two-dimensional face $F \preceq \Delta$ whose cone $C(F)$ is singular,
$\ell(F^\circ)$ points at which the germ of $X$ is analytically isomorphic to
$C(F)$, where $\ell(F^\circ)$ is the lattice length of the edge
$F^\circ \preceq \Delta^\circ$ dual to $F$.  The Altmann--Corti--Filip--Petracci
trichotomy, recalled as \cite[Thm.~2.2]{paper1},
classifies $C(F)$ by the combinatorics of $F$ into rigid \textup{(R)},
deformable-but-non-smoothable \textup{(D)}, and smoothable \textup{(S)}; a
single face of the first two types forces $X$ to admit no smoothing, and
\cite{paper2} studies when the type-\textup{(S)} configurations smooth.

Mirror symmetry, in the Batyrev picture, is the polar duality
$\Delta \leftrightarrow \Delta^\circ$, which for $4$-polytopes exchanges
two-dimensional faces with edges.  This note asks the mirror question: when
are $X$ and the mirror $X^\circ \subset \PP_{\Delta^\circ}$ \emph{both}
isolated-singular (we then call the pair \emph{doubly isolated}), and
which local models, in particular which
non-smoothable ones, can occur in that situation?

The entry point is exact and elementary.  For a two-dimensional face $F$
write $u_1, u_2 \in \MM$ for the primitive inner normals of the two facets of
$\Delta$ containing $F$, so $F^\circ = \conv\{u_1,u_2\}$, and let
$\MM_F = \MM \cap (\RR u_1 + \RR u_2)$.  We call the two-dimensional cone
$\Cone(u_1,u_2)$ the \emph{normal cone} of $F$ and its dual the
\emph{transverse cone} of $\Delta$ along $F$ (the local model of
$\Delta$ in the directions normal to $F$), and we say $\Delta$ is
\emph{transversally smooth} along $F$ when this cone is smooth.

\begin{theorem}[= Lemma~\ref{lem:mult} with Remark~\ref{rem:transverse}]\label{thm:intromult}
$\ell(F^\circ) = [\,\MM_F : \ZZ u_1 + \ZZ u_2\,]$, the multiplicity of the
normal cone $\Cone(u_1,u_2)$ of $F$ (equivalently, of its dual, the
transverse cone of $\Delta$ along $F$).  In particular $\ell(F^\circ)=1$ if
and only if the transverse cone is smooth.
\end{theorem}

$X$ is isolated-singular exactly when $\Delta$ has unit edges, and
$X^\circ$ exactly when every two-dimensional face of $\Delta$ has
$\ell(F^\circ)=1$; a face with $\ell(F^\circ) = m \ge 2$ contributes to
$X^\circ$ a curve of transverse $A_{m-1}$ singularities
(Proposition~\ref{prop:corr}(2)).  We obtain:

\begin{corollary}\label{cor:introboth}
$X$ and $X^\circ$ are both isolated-singular if and only if $\Delta$ has unit
edges and every two-dimensional face of $\Delta$ is transversally smooth.  We
call such $\Delta$ \emph{both-sides unit}.  The condition is self-dual:
$\Delta$ is both-sides unit if and only if $\Delta^\circ$ is.
\end{corollary}

Our main results determine the both-sides unit polytopes, and the local
models they carry, \emph{completely}, subject to one explicit input
assumption.  We use the following \emph{database hypothesis}: the
per-vertex-count files described in \S\ref{sec:census} contain, without
repetition, exactly the members of the Kreuzer--Skarke classification in
their stated vertex range; the unique remaining $36$-vertex member is
included separately.  Under this hypothesis the scan covers all
$473{,}800{,}776$ reflexive $4$-polytopes \cite{KreuzerSkarke00}.  Every
polytope returned by the scan is then re-verified in exact arithmetic, so
the classification statements below have finite computer-assisted proofs
conditional only on that database hypothesis.  The transverse identity
and the assertions about each explicitly displayed polytope are
unconditional.

\begin{theorem}[classification of both-sides local models]\label{thm:class}
Under the database hypothesis above, there are exactly $590$ both-sides
unit reflexive $4$-polytopes.  Every
two-dimensional face of every one of them is
\begin{itemize}
\item a triangle (the germ is a cyclic quotient point, or smooth), or
\item a zonotope (a centrally-symmetric unit-edge polygon, always of
type \textup{(S)}), or
\item a reflexive polygon: one of the five unit-edge polygons with exactly
one interior lattice point, i.e.\ the Fano polygons of the five
smooth toric del Pezzo surfaces.
\end{itemize}
All five reflexive polygons occur.  The complete face inventory is given in
Table~\textup{\ref{tab:occ}}.
\end{theorem}

\begin{theorem}[the non-smoothable germs of doubly isolated pairs]
\label{thm:dichot}
Under the database hypothesis, exactly three of the $590$ both-sides unit
polytopes carry a non-smoothable
two-dimensional face.  Writing $[k,i]$ for a unit-edge polygon with $k$
edges and $i$ interior lattice points, the germs occurring are:
\begin{itemize}
\item $\frac13(1,1,1)$ (the cone over the $[3,1]$ triangle), six times on
one $9$-vertex polytope;
\item $\frac15(1,1,3)$ (the cone over the $[3,2]$ triangle), ten times on
a single simplex;
\item the anticanonical cone over $\FF_1$, exactly \emph{once}, on a
unique $22$-vertex polytope.
\end{itemize}
In particular no deformable-but-non-smoothable
\textup{(D)} face occurs on any both-sides unit polytope.
\end{theorem}

\begin{corollary}\label{cor:defonly}
Assume the database hypothesis.  Let $X$ be a unit-edge Batyrev
Calabi--Yau threefold with a singular point of
type \textup{(D)}, for instance the pentagon-cone threefolds of
\cite[Thms.~5.3, 5.8]{paper1}.  Then the Batyrev mirror $X^\circ$ is
non-isolated: it carries a curve of transverse $A_{m-1}$ singularities for
some $m \ge 2$.  The
deformable-but-non-smoothable phenomenon never occurs on a doubly isolated
mirror pair.
\end{corollary}

The one non-quotient exception is the most interesting object of this note.

\begin{theorem}[the unique $\FF_1$ mirror pair]\label{thm:pair}
Under the database hypothesis, up to lattice isomorphism there is exactly
one doubly isolated Batyrev mirror
pair whose singularities are not all quotient points or smoothable: the pair
$(\Delta_{\FF_1}, \Delta_{\FF_1}^\circ)$ of Theorem~\textup{\ref{thm:pairbody}},
with $22$ and $26$ vertices.  The generic anticanonical
$X \subset \PP_{\Delta_{\FF_1}}$ has $20$ ordinary double points, one
$\mathrm{dP}_6$-cone point, and one point whose germ is the anticanonical
cone over $\FF_1$; since that germ admits no smoothing at all, $X$
\emph{admits no smoothing} \cite[Cor.~4.3]{paper1}.  Its mirror $X^\circ$
has $26$ ordinary double
points, two $\mathrm{dP}_7$-cone points and two $\mathrm{dP}_6$-cone
points; every singular germ of $X^\circ$ is of type \textup{(S)}.  The maximal
projective crepant partial resolutions have
$(h^{1,1},h^{2,1}) = (20,26)$ and $(26,20)$ respectively.
\end{theorem}

Thus the Batyrev mirror of a Calabi--Yau threefold that is \emph{frozen}
(non-smoothable, carrying the cone over one of the two surfaces excluded in
Gross's smoothing theorem for
primitive contractions \cite[Thm.~5.8]{Gross97}) can be isolated-singular
with every singular point locally smoothable.  Notably, every dual edge on
the mirror side has length one, so $X^\circ$ lies entirely in the
$\ell(F^\circ)=1$ regime of \cite{paper2}.  Its $26$ nodes pose the
\emph{cross-face} problem of \cite[\S7]{paper2}.  Exact row reduction of
their $26$ diagonal-relation vectors gives rank $18$ and two coloops, so
the node subsystem has no relation with every coefficient non-zero
(Remark~\ref{rem:node-subsystem}).  This would obstruct smoothing in an
all-nodal threefold with a projective small resolution.  The actual
$X^\circ$, however, also has four divisorial del Pezzo-cone points, so
Friedman's nodal criterion does not apply to the mixed singularity
configuration: their local deformation directions may participate in the
global obstruction map.  Its del Pezzo-cone germs are precisely those that
\cite[\S6]{paper2} places beyond Gross's theorems and for which no
Friedman-type obstruction is known.  Thus any smoothing would have to be
genuinely mixed, and whether one exists remains open
(Question~\ref{q:mirror-smooth}).

Theorems~\ref{thm:class}, \ref{thm:dichot} and~\ref{thm:pair} answer the
mirror question of
\cite[\S6.3]{paper1} for the non-smoothable local models: among them,
only the cyclic quotients and the $\FF_1$-cone,
including the germs of Gross's two exceptional surfaces $\PP^2$ and
$\FF_1$ \cite[Thm.~5.8]{Gross97}, admit doubly isolated mirror
realizations, and the $\FF_1$-cone does so exactly once in the entire toric
hypersurface landscape.

The opposite extreme of the classification carries a matching uniqueness.
Exactly eight of the $590$ polytopes have all two-dimensional faces smooth,
so that $X$ itself is smooth; for exactly one of them the mirror is smooth
as well: the self-dual $24$-cell, whose hypersurface is the classical
self-mirror threefold with $(h^{1,1},h^{2,1}) = (20,20)$ \cite{Braun12}.
The smooth Batyrev mirror pair is thus unique in the entire
classification (Example~\ref{ex:24cell}).

\subsection*{A caution about low-vertex evidence}
An earlier version of this note conjectured that every two-dimensional face
of a both-sides unit polytope is a triangle or a zonotope; equivalently,
that no \emph{asymmetric} face (one that is neither a triangle nor
centrally symmetric) ever occurs.  The conjecture is \emph{true} for all
$395{,}406{,}329$ reflexive
$4$-polytopes with at most $18$ vertices ($83.5\%$ of the classification)
and \emph{false} from $19$ vertices on, exactly where both-sides unit
polytopes first
become common (Table~\ref{tab:dist}): the $\mathrm{dP}_7$ pentagon occurs
($87$ faces, on polytopes with $19$ to $31$ vertices), as does the $\FF_1$
quadrilateral (once).  See \S\ref{sec:refuted} for the refuted statements.
Low-vertex evidence, however extensive, samples a biased corner of the
landscape.

All mathematical predicates and reported invariants in the scan, the exact
re-verification of all $590$ polytopes, and the mirror-pair calculation use
integer or rational arithmetic; the sole floating-point ordering operation
is disclosed in \S\ref{subsec:method}.  The programs build on those of
\cite{paper1} and are included as ancillary files;
\texttt{both\_sides\_census.py} re-derives every number and
theorem in this note from the scan results.  The
code and the scan data are public at
\url{https://github.com/bwuebben/calabi-yau-smoothability}, where every
count in this note is recomputable from
\texttt{src/both\_sides\_census.py}.

\subsection*{Acknowledgements}
This note completes the series \cite{paper1,paper2}; the question of
smoothing canonical Calabi--Yau singularities was posed to the author by
Mark Gross many years ago.

\section{Preliminaries}\label{sec:prelim}

We keep the notation of \cite{paper1}.  For a lattice polygon $Q$ with unit
edges placed at height one, $C(Q)$ denotes the associated Gorenstein toric
affine threefold; $Q$ is of type \textup{(R)}, \textup{(D)} or \textup{(S)}
according as its edge-vector multiset has no proper zero-sum sub-multiset,
has one but no partition into unit segments and unimodular triples, or has
such a partition \cite{Altmann97,CFP}.  Type \textup{(S)} polygons are
exactly those with smoothable cone; a single germ of type \textup{(R)} or
\textup{(D)}, one whose miniversal base has no smoothing component,
obstructs every global smoothing of $X$
\cite[Lemma~4.1, Cor.~4.3]{paper1}.  We write $[k,i]$ for the class of a
unit-edge polygon with $k$ edges and $i$ interior lattice points.  A
unit-edge polygon is a
\emph{zonotope} (a Minkowski sum of unit segments) if and only if it is
centrally symmetric, and every zonotope is of type \textup{(S)}.  A unit-edge
polygon with exactly one interior lattice point is reflexive.  Indeed,
translate that point to the origin and let the polygon have $k$ edges.
Pick's theorem gives normalized area $k$, while the fan from the origin
decomposes it into $k$ lattice triangles of positive integral normalized
areas, one over each unit edge.  Every one therefore has normalized area
one, so the origin has lattice distance one from every edge.  There are
five, the Fano polygons of $\PP^2$, $\FF_1$, $\PP^1{\times}\PP^1$,
$\mathrm{dP}_7$, $\mathrm{dP}_6$ (the subscript is the degree:
$\mathrm{dP}_7 = \mathrm{Bl}_2\PP^2$, $\mathrm{dP}_6 = \mathrm{Bl}_3\PP^2$),
the cone $C(Q)$ being the anticanonical
cone over the corresponding surface.  Of the five, the $\PP^2$ triangle and
the $\FF_1$ quadrilateral are rigid and the other three are of type
\textup{(S)} \cite[Rem.~2.6]{paper1}.  For a reflexive $4$-polytope $\Delta$ we write
$\Delta^\circ = \{u : \langle u,v\rangle \ge -1\ \forall v \in \Delta\}$; a
face $G \preceq \Delta$ has dual face $G^\circ \preceq \Delta^\circ$ with
$\dim G + \dim G^\circ = 3$.  By
\cite[Prop.~3.1]{paper1}, for $\Delta$ with unit edges the generic
$X \in |{-K}|$ has, over each two-dimensional face $F$ with $C(F)$
singular, exactly
$\ell(F^\circ)$ singular points with germ $C(F)$, and no other singular
points; and an edge of $\Delta$ of lattice length $m \ge 2$ produces a curve
of transverse $A_{m-1}$ singularities \cite[Rem.~3.2]{paper1}.  Throughout,
$\Delta \subset \NN_\RR$ is the \emph{fan} polytope, Batyrev's
$\Delta^\ast$, so our $\Delta^\circ$ is his $\Delta$; the same
convention is used in \cite{paper1,paper2}.

\section{The transverse-multiplicity identity}\label{sec:mult}

Let $F$ be a codimension-two face of a reflexive polytope
$\Delta \subset \NN_\RR$ of arbitrary dimension, and let
$u_1, u_2 \in \MM$ be the primitive inner
normals of the two facets of $\Delta$ containing $F$, so that
$\langle u_i, \cdot\rangle = -1$ on $F$ and $F^\circ = \conv\{u_1, u_2\}$.
\footnote{The $u_i$ are vertices of $\Delta^\circ$; \cite{paper1} and
the ancillary programs list the sign-flipped functionals taking the value
$+1$ on the facet.  Lattice lengths are unaffected.}
Put $\MM_F = \MM \cap (\RR u_1 + \RR u_2)$, a rank-two saturated sublattice
of $\MM$ containing $u_1, u_2$ as primitive vectors.

\begin{lemma}\label{lem:mult}
$\ell(F^\circ) = [\,\MM_F : \ZZ u_1 + \ZZ u_2\,]$.
\end{lemma}

\begin{proof}
Both sides depend only on the pair $(u_1,u_2)$ inside the rank-two lattice
$\MM_F$.  Choose a $\ZZ$-basis of $\MM_F$ in which $u_1 = (0,-1)$; this is
possible because $u_1$ is primitive.  Write $u_2 = (p,q)$; here $p \neq 0$, since $p = 0$ would force
$u_2 = \pm u_1$, impossible as $F^\circ = \conv\{u_1,u_2\}$ is an edge of
$\Delta^\circ$ and $0$ is interior to $\Delta^\circ$.  Then $p =
\det(u_1,u_2) = \pm[\,\MM_F : \ZZ u_1 + \ZZ u_2\,]$, and $\ell(F^\circ) =
\ell\bigl(\conv\{u_1,u_2\}\bigr) = \gcd(u_2 - u_1) = \gcd(p,\,q+1)$,
the two gcds agreeing because $\MM_F$ is saturated in $\MM$, so lattice
length in $\MM$ may be computed in the chosen basis of $\MM_F$.

Let $v$ be a vertex of $F$ and let $\bar v \in \NN_F := \NN/(\NN \cap
u_1^\perp \cap u_2^\perp)$ be its image; since $u_1,u_2$ vanish on
$\NN \cap u_1^\perp \cap u_2^\perp$ they descend to elements of
$\MM_F = \NN_F^\vee$, and $\bar v$ is a lattice point of $\NN_F \cong \ZZ^2$
with $\langle u_i, \bar v\rangle = -1$.  Writing $\bar v = (a,b)$ in the
basis of $\NN_F$ dual to the chosen basis of $\MM_F$,
$\langle u_1, \bar v\rangle = -b = -1$ gives $b = 1$, and
$\langle u_2, \bar v\rangle = pa + q = -1$ gives $p a = -1-q$, so
$p \mid (q+1)$.  Hence $\gcd(p, q+1) = |p|$, and
$\ell(F^\circ) = |p| = [\,\MM_F : \ZZ u_1 + \ZZ u_2\,]$.
\end{proof}

\begin{remark}\label{rem:transverse}
Lemma~\ref{lem:mult} is elementary and may be known; what matters here is
the reading it provides.  The proof works in every dimension for a
codimension-two face.  The right-hand side is the multiplicity of the
two-dimensional cone $\Cone(u_1,u_2)$, which is the normal cone of $F$ and
dually the transverse cone of $\Delta$ along $F$ (for two-dimensional cones
the multiplicities of a cone and its dual agree).  Thus \emph{the dual-edge
length counts exactly the transverse singularity multiplicity}, and
$\ell(F^\circ)=1$ is the statement that $\Delta$ is transversally smooth
along $F$.  (Transversal smoothness concerns only the directions normal to
$F$: the germ $C(F)$ itself may well be singular.)  The lattice-point
hypothesis on $v$ is essential: for an
arbitrary pair of primitive vectors the two sides can differ.
\end{remark}

Corollary~\ref{cor:introboth} is immediate: $X$ is isolated-singular iff
$\Delta$ has unit edges; $X^\circ$ is isolated-singular iff $\Delta^\circ$
has unit edges, i.e.\ every edge $F^\circ$ of $\Delta^\circ$ has
$\ell(F^\circ)=1$, i.e.\ by Lemma~\ref{lem:mult} every two-dimensional face
of $\Delta$ is transversally smooth.  Self-duality follows since the
conditions ``unit edges'' and ``all dual edges of length one'' are exchanged
by $\Delta \leftrightarrow \Delta^\circ$.

\section{The mirror correspondence of the singular strata}\label{sec:corr}

\begin{proposition}\label{prop:corr}
Let $\Delta$ be reflexive with unit edges and $F$ any two-dimensional
face.  Under polar duality $F$ corresponds to the edge
$F^\circ$ of $\Delta^\circ$, and:
\begin{enumerate}
\item if $C(F)$ is singular, $X$ has $\ell(F^\circ)$ isolated points of
type $C(F)$ over $F$ (and none over $F$ otherwise);
\item $X^\circ$ has, along the surface
$V(\sigma_{F^\circ}) \subset \PP_{\Delta^\circ}$ (the orbit closure of
the cone $\sigma_{F^\circ}$ over the edge $F^\circ$ in the face fan of
$\Delta^\circ$), a curve of transverse $A_{\ell(F^\circ)-1}$
singularities (and none if $\ell(F^\circ)=1$).
\end{enumerate}
In particular $X^\circ$ is isolated-singular if and only if
$\max_F \ell(F^\circ) = 1$, i.e.\ $\Delta$ is both-sides unit.
\end{proposition}

\begin{proof}
(1) is \cite[Prop.~3.1]{paper1}.  For (2), the edge $F^\circ$ of
$\Delta^\circ$ has lattice length $\ell(F^\circ)$; an edge of lattice length
$m$ produces on the generic anticanonical hypersurface a curve of transverse
$A_{m-1}$ singularities (\S\ref{sec:prelim}), and none when $m = 1$.  Hence
$X^\circ$ is isolated-singular iff every two-dimensional face of $\Delta$
has $\ell(F^\circ) = 1$.
\end{proof}

Thus the $\ell(F^\circ)$ singular points of $X$ over $F$ and the
$A_{\ell(F^\circ)-1}$-curve of $X^\circ$ correspond under polar duality,
and the two
threefolds are simultaneously isolated-singular exactly on the both-sides
unit polytopes.  Whenever $\Delta$ has \emph{any} face with
$\ell(F^\circ) \ge 2$ (the generic situation), the mirror is
non-isolated and its singularities leave the isolated framework of
\cite{Altmann97,CFP}.  Filip \cite{Filip25} constructs, from a Laurent
polynomial, a distinguished formal deformation of the associated
non-isolated Gorenstein toric pair and proves that its general fibre is
smooth exactly when the polynomial is $0$-mutable; this does not classify
all components of every non-isolated germ.

\section{Enumeration over the Kreuzer--Skarke
classification}\label{sec:census}

\subsection{Methodology}\label{subsec:method}
The both-sides unit condition is a condition on faces: all edges of
lattice length one
and, by Lemma~\ref{lem:mult}, all dual edges of length one.  We tested it
over
the per-vertex-count copy of the Kreuzer--Skarke classification used in
\cite[\S6.1]{paper1}, with the production implementation built for that
scan.  Subject to the database hypothesis in the introduction, this covers
the complete classification.  All facet incidences, lattice lengths,
determinants, and final predicates are checked in exact integer arithmetic,
with no heuristic pruning.  The implementation uses floating-point
\texttt{atan2} only to choose a cyclic ordering of the vertices of an
already identified planar face; the ensuing edge and lattice tests are
integral.  Implementation correctness is controlled in three complementary ways.
(i) Before scanning each file of the database, the implementation was
re-checked against an
independently written reference implementation on random samples from that
file.
(ii) The $3{,}905{,}789$ polytopes with at most nine vertices were scanned
in
full by both implementations, with identical results.  (iii) The $590$
polytopes found pass a global consistency check under polar duality: the
both-sides condition
is self-dual (Corollary~\ref{cor:introboth}), and the polar of every
polytope found,
computed exactly, is again both-sides unit and agrees with one of the
$590$ in vertex
count, facet count and complete face inventory, an invariance that
would generically fail for a false
negative whose polar partner was detected.  An invariant collision
(another member of the $590$ sharing its vertex count, facet count and
face inventory) could in principle mask a miss, so this control is
supporting evidence beside (i) and (ii), not an independent proof.  The one polytope
missing from the copy of the database we used, which is organized by
vertex count, is the $36$-vertex hexagon product (see
\cite[\S6.1]{paper1}); it is included.
Every polytope passing the scan is then re-verified, and its faces
classified, by the reference implementation; the
ancillary scripts reproduce the full scan from the public database.  The
scan results are included as ancillary data files, and the program
\texttt{both\_sides\_census.py} re-derives every statement of
Theorems~\ref{thm:class}, \ref{thm:dichot} and~\ref{thm:pair} from them.

\subsection{The classification}
\begin{table}[ht]
\adjustbox{max width=\textwidth}{%
\begin{tabular}{lrrrrrrrrrrrrrrrrrrrrrrrrr}
\toprule
vertices & 5 & 6 & 8 & 9 & 10 & 12 & 15 & 16 & 17 & 18 & 19 & 20 & 21 & 22
& 23 & 24 & 25 & 26 & 27 & 28 & 29 & 30 & 31 & 33 & 36\\
\midrule
\# & 2 & 1 & 1 & 1 & 1 & 1 & 1 & 4 & 4 & 7 & 12 & 24 & 37 & 69 & 84 & 98 &
85 & 64 & 42 & 24 & 13 & 9 & 4 & 1 & 1\\
\bottomrule
\end{tabular}}
\medskip
\caption{The $590$ both-sides unit reflexive $4$-polytopes by vertex count
(no others occur; in particular none with $7$, $11$, $13$, $14$, $32$, $34$
or $35$ vertices).  The population is closed under polar duality
(Corollary~\textup{\ref{cor:introboth}}) and peaks at $24$
vertices.}\label{tab:dist}
\end{table}

\begin{table}[ht]
\begin{center}
\begin{tabular}{llrl}
\toprule
face $[k,i]$ & kind & count & germ of $X$ over it \\
\midrule
$[3,0]$ & triangle & $25{,}873$ & smooth point \\
$[3,1]$ & triangle & $6$ & $\frac13(1,1,1)$ \textbf{(rigid)} \\
$[3,2]$ & triangle & $10$ & $\frac15(1,1,3)$ \textbf{(rigid)} \\
$[4,0]$ & zonotope & $15{,}652$ & ordinary double point \\
$[4,1]$ & zonotope & $99$ & cone over $\PP^1{\times}\PP^1$ \\
$[4,1]$ & \emph{asymmetric} & $1$ & cone over $\FF_1$ \textbf{(rigid)} \\
$[4,2]$ & zonotope & $9$ & type (S) \\
$[5,1]$ & \emph{asymmetric} & $87$ & cone over $\mathrm{dP}_7$ \\
$[6,1]$ & zonotope & $395$ & cone over $\mathrm{dP}_6$, $2$ smoothing components \\
$[6,2]$ & zonotope & $2$ & type (S) \\
$[6,7]$ & zonotope & $3$ & type (S) \\
$[8,4]$ & zonotope & $1$ & type (S), $3$ smoothing components \\
\bottomrule
\end{tabular}
\end{center}
\medskip
\caption{Every two-dimensional face of every both-sides unit polytope, by
class ($k$ = edges, $i$ = interior points).  ``Asymmetric'' = neither a
triangle nor centrally symmetric; only the two non-centrally-symmetric
reflexive polygons occur, and each germ appears on $X$ exactly once per
face since all dual edges have length one.  Type-\textup{(S)} germs have
one smoothing component except where noted.}\label{tab:occ}
\end{table}

\begin{proof}[Proof of Theorems~\textup{\ref{thm:class}}
and~\textup{\ref{thm:dichot}}]
Both theorems are exhaustive verifications, summarized in
Tables~\ref{tab:dist} and~\ref{tab:occ}; the completeness of the
underlying enumeration is the content of \S\ref{subsec:method}.
Granting it, the scan of all $473{,}800{,}776$ polytopes yielded
$590$ both-sides unit ones; exact re-analysis of each classifies all their
two-dimensional faces into the twelve classes of Table~\ref{tab:occ}, each
of which is a triangle, a zonotope, or one of the five reflexive polygons
(the $[4,1]$ zonotope is the $\PP^1{\times}\PP^1$ diamond, the $[6,1]$
zonotope the $\mathrm{dP}_6$ hexagon, the $[3,1]$ triangle the $\PP^2$
triangle; the asymmetric $[5,1]$ and $[4,1]$ classes are the $\mathrm{dP}_7$
pentagon and the $\FF_1$ quadrilateral, identified by explicit
$GL_2(\ZZ)$-equivalences).  The non-smoothable classes occurring are the
rigid triangles $[3,1]$, $[3,2]$ and the $\FF_1$ quadrilateral; the
$[3,2]$ faces are the ten faces of the quotient simplex of
Example~\ref{ex:pair}, the $[3,1]$ faces lie on the $9$-vertex polytope of
Example~\ref{ex:pair}, and
the $\FF_1$ face on the unique polytope of
Theorem~\ref{thm:pairbody}.  No face of type \textup{(D)} occurs.
\end{proof}

\begin{proof}[Proof of Corollary~\textup{\ref{cor:defonly}}]
A type-\textup{(D)} point of $X$ is the cone over a type-\textup{(D)} face
$F$.  By Theorem~\ref{thm:dichot}, $\Delta$ is then not both-sides unit, so
some face $F'$ has $\ell(F'^\circ) \ge 2$ and Proposition~\ref{prop:corr}(2)
gives a singular curve on $X^\circ$.
\end{proof}

\subsection{The refuted conjectures}\label{sec:refuted}
The first version of this note conjectured that every two-dimensional face
of a both-sides unit polytope is a triangle or a zonotope, and reduced this
to a local ``balancing'' claim: an asymmetric face $F$ with
$\ell(F^\circ)=1$, all faces of whose star have unit edges, forces some
face of its star to have $\ell \ge 2$.  The local claim
had survived a complete check of the $38{,}571$ candidate configurations
arising among the $1{,}037{,}834$ polytopes with at most eight vertices.
Table~\ref{tab:occ} nevertheless
refutes both statements: \mbox{$88$ asymmetric faces occur, on $55$ polytopes}.
The smallest
counterexample has $19$ vertices:
\[
\begin{gathered}
\Delta_{19}=\conv\{(1,0,0,0),\,(0,1,0,0),\,(0,0,1,0),\,(0,0,-1,0),\,
(0,-1,0,0),\\
(1,-1,-1,0),\,(0,0,0,1),\,(-1,1,1,-1),\,(0,-1,0,1),\,(0,0,-1,1),\\
(-1,1,0,-1),\,(-1,0,1,-1),\,(-1,0,1,0),\,(-1,1,0,0),\,(0,-1,-1,0),\\
(0,-1,-1,1),\,(-1,0,0,-1),\,(-2,1,1,-1),\,(-1,0,0,1)\},
\end{gathered}
\]
a both-sides unit polytope one of whose faces is the $\mathrm{dP}_7$
pentagon; its hypersurface is a Calabi--Yau threefold with $14$ nodes and
one $\mathrm{dP}_7$-cone point, all with isolated-singular mirror.  That face
is
\[
P = \conv\bigl\{(0,-1,-1,0),\ (0,-1,-1,1),\ (-1,0,0,-1),\
(-2,1,1,-1),\ (-1,0,0,1)\bigr\},
\]
the fifteenth to nineteenth entries of the list, lying in the two facets with
inner normals $u_1 = (1,0,1,0)$ and $u_2 = (1,1,0,0)$; their difference
$(0,-1,1,0)$ is primitive, so $\ell(P^\circ) = 1$.  In the induced lattice
its edge multiset is $\{(0,-1),\,(1,-1),\,(1,0),\,(-1,2),\,(-1,0)\}$, carried
by
$M=\left(\begin{smallmatrix}-1&0\\1&1\end{smallmatrix}\right) \in GL_2(\ZZ)$
onto $\{(0,-1),\,(1,-1),\,(1,1),\,(-1,1),\,(-1,0)\}$, the
$\mathrm{dP}_7$ pentagon of \cite[Rem.~2.6]{paper1}: the unique
unit-edge class with
five vertices and one interior lattice point.  The
conjecture is \emph{true} on the $395{,}406{,}329$ polytopes with at most
$18$ vertices and false from $19$ on --- precisely where both-sides unit
polytopes stop being sporadic (Table~\ref{tab:dist}).  We let the corrected
statement stand as Theorem~\ref{thm:class}, and the conceptual question it
raises as Question~\ref{q:why}.

\begin{remark}[why not ``simplicial'']\label{rem:selfdual}
The stronger guess that a both-sides unit polytope is simplicial fails for a
structural reason: the condition is self-dual, and ``simplicial and simple''
forces a simplex, contradicting the existence of both-sides unit polytopes
with more than five vertices.  The non-triangular faces are genuine; the
content of Theorem~\ref{thm:class} is that they are zonotopes or reflexive.
\end{remark}

\begin{example}[the low-vertex stratum]\label{ex:pair}
Among the $3{,}905{,}789$ polytopes with at most nine vertices exactly five
are both-sides unit; four of them form two polar pairs.  One pair is the
smooth simplex and the simplex
\[
\conv\{(1,0,0,0),\,(0,1,0,0),\,(2,4,5,0),\,(3,3,0,5),\,(-6,-8,-5,-5)\},
\]
whose ten
faces are the $[3,2]$ rigid triangles of Table~\ref{tab:occ} (quotient
points of type $\frac15(1,1,3)$) and whose polar is smooth.  A second pair consists of a
simplicial six-vertex polytope with all faces smooth and its nine-vertex
polar, carrying the six $\frac13(1,1,1)$ triangles and nine $[4,2]$
zonotopes.  In each pair, the non-smoothable member mirrors a \emph{smooth}
threefold.  The fifth polytope has eight vertices and sixteen facets, and
all thirty-two of its faces are smooth triangles, so its hypersurface $X$
is smooth.  It shows that closedness under polar duality
(Corollary~\ref{cor:introboth}) does not respect bounds on the number of
vertices: its
polar is a sixteen-vertex both-sides unit polytope all twenty-four of whose
faces are $[4,1]$ diamonds, so the mirror $X^\circ$ carries twenty-four
cone-over-$\PP^1{\times}\PP^1$ points.
\end{example}

\begin{example}[the top of the classification]\label{ex:pair36}
The $36$-vertex product $H \times H$ of two reflexive hexagons is the
reflexive $4$-polytope with the most vertices and the one polytope absent
from the copy of the Kreuzer--Skarke database used in \cite{paper1}.  It is
both-sides unit; its faces are $36$ unit squares and $12$ $\mathrm{dP}_6$
hexagons, so its hypersurface has $48$ singular points, all of type
\textup{(S)}, while its polar, the $12$-vertex free sum
$H^\circ \oplus H^\circ$, has $72$ smooth triangular faces and smooth
hypersurface.
\end{example}

\begin{example}[the unique smooth mirror pair]\label{ex:24cell}
Exactly eight of the $590$ polytopes have every two-dimensional face
smooth, so that $X$ itself is smooth.  For exactly one of them is the polar
of the same kind: the $24$-cell, self-dual, with $24$ octahedral facets,
$96$ unimodular triangular faces, and no lattice points besides the origin
and its $24$ vertices.  Its generic anticanonical hypersurface is the
classical self-mirror Calabi--Yau threefold with
$(h^{1,1},h^{2,1}) = (20,20)$ and $\chi = 0$; the $(1,1)$-threefolds of
Braun \cite{Braun12} are free quotients of special symmetric members of
this family.  The scan shows the uniqueness:
among all $473{,}800{,}776$ reflexive $4$-polytopes, the $24$-cell is the
only one whose generic anticanonical hypersurface and mirror are
\emph{both} smooth.
\end{example}

\section{The unique $\FF_1$ mirror pair}\label{sec:pair}

\begin{theorem}\label{thm:pairbody}
Let
\[
\begin{gathered}
\Delta_{\FF_1}=\conv\{(1,0,0,0),\,(0,1,0,0),\,(1,-1,0,0),\,(0,0,1,0),\,
(0,0,0,1),\,(0,0,1,-1),\\
(0,0,-1,1),\,(0,0,0,-1),\,(0,0,-1,0),\,(-1,1,0,0),\,(0,-1,0,0),\,
(-1,1,-1,1),\\
(0,-1,0,-1),\,(0,-1,-1,0),\,(-1,1,-1,0),\,(-1,0,0,-1),\,(-1,0,-1,1),\\
(-1,-1,0,-1),\,(-2,1,-1,1),\,(-2,1,-1,0),\,(-1,-1,-1,0),\,(-2,0,-1,0)\}
\subset \ZZ^4 .
\end{gathered}
\]
Then $\Delta_{\FF_1}$ is a both-sides unit reflexive polytope with $22$
vertices and $26$ facets, and:
\begin{enumerate}
\item its $72$ two-dimensional faces are $50$ smooth triangles, $20$ unit
squares, one $\mathrm{dP}_6$ hexagon, and one quadrilateral
$GL_2(\ZZ)$-equivalent to the $\FF_1$ polygon; every dual edge has lattice
length one;
\item the generic anticanonical $X \subset \PP_{\Delta_{\FF_1}}$ is a
Calabi--Yau threefold whose singular locus consists of $20$ ordinary double
points, one $\mathrm{dP}_6$-cone point, and one point with germ the
anticanonical cone over $\FF_1$; since one germ without a smoothing
component obstructs every global smoothing (\S\ref{sec:prelim}), $X$
admits no smoothing \cite[Cor.~4.3]{paper1};
\item the polar $\Delta_{\FF_1}^\circ$ ($26$ vertices, $22$ facets) is
both-sides unit with $68$ two-dimensional faces: $38$ smooth triangles, $26$ unit squares,
two $\mathrm{dP}_7$ pentagons and two $\mathrm{dP}_6$ hexagons; the mirror
$X^\circ$ has $26$ ordinary double points, two $\mathrm{dP}_7$-cone points
and two $\mathrm{dP}_6$-cone points, each with a local smoothing;
\item the maximal projective crepant partial resolutions satisfy
$(h^{1,1},h^{2,1})(\widehat X) = (20,26)$ and
$(h^{1,1},h^{2,1})(\widehat{X^\circ}) = (26,20)$;
\item under the database hypothesis, $\Delta_{\FF_1}$ is the \emph{only}
both-sides unit reflexive
$4$-polytope, up to lattice isomorphism, carrying a
non-smoothable face that is not a triangle.
\end{enumerate}
\end{theorem}

\begin{proof}
(1)--(4) are finite checks on the explicitly listed polytope, a face
enumeration reproducible in any lattice-polytope package, carried out
in exact arithmetic by the ancillary files.\footnote{Ancillary files
\texttt{both\_sides\_census.py} and \texttt{face\_data.py}.}  We record
the $\FF_1$ face itself, since it is the object the statement is about.  Its
four vertices are the fourth, eighteenth, nineteenth and twenty-second entries
of the list above,
\[
F \;=\; \conv\bigl\{(0,0,1,0),\ (-1,-1,0,-1),\ (-2,1,-1,1),\
(-2,0,-1,0)\bigr\},
\]
and the two facets of $\Delta_{\FF_1}$ containing $F$ have inner normals
\[
u_1 = (1,1,-1,-1), \qquad u_2 = (1,0,-1,0),
\]
in the normalization of \S\ref{sec:intro}: both take the value $-1$ at each
of the four vertices of $F$ and are $\ge -1$ on $\Delta_{\FF_1}$.  Their
difference $u_1 - u_2 = (0,1,0,-1)$ is primitive, so
$\ell(F^\circ) = 1$.  In the induced lattice on the affine span of $F$ the
edge multiset is
$\{(0,-1),\,(1,-1),\,(1,1),\,(-2,1)\}$, and
\[
M=\begin{pmatrix}-1&-1\\-1&0\end{pmatrix} \in GL_2(\ZZ),
\qquad \det M = -1,
\]
carries it onto $E(Q_A) = \{(-2,-1),\,(0,-1),\,(1,0),\,(1,2)\}$, the edge
multiset of the $\FF_1$ polygon of \cite[Rem.~2.6]{paper1}; so $C(F)$ is the
anticanonical cone over $\FF_1$.

The $\mathrm{dP}_6$ hexagon of (1) is recorded the same way.  It is the face
\[
\begin{gathered}
G = \conv\bigl\{(0,0,-1,0),\ (0,-1,-1,0),\ (-1,1,-1,0),\\
(-2,1,-1,0),\ (-1,-1,-1,0),\ (-2,0,-1,0)\bigr\}
\end{gathered}
\]
(the ninth, fourteenth, fifteenth, twentieth, twenty-first and
twenty-second entries of the list), lying in the two facets with inner
normals $u_1 = (0,0,1,0)$ and $u_2 = (0,0,1,1)$, whose difference
$(0,0,0,-1)$ is primitive, so $\ell(G^\circ) = 1$.  Its edge multiset in the
induced lattice is already the standard hexagon,
\[
E(G) = \{(0,-1),\,(1,-1),\,(1,0),\,(0,1),\,(-1,1),\,(-1,0)\},
\]
so no change of basis is needed: $C(G)$ is the anticanonical cone over
$\mathrm{dP}_6$, the class with two smoothing components in the catalogue of
\cite[Rem.~2.6]{paper1}.  The Hodge numbers are
given by Batyrev's formulas \cite[\S4]{Batyrev94}, evaluated in exact
arithmetic.  The mirror swap in (4) is forced by Batyrev duality and serves as
a consistency check.  Conditional on the database hypothesis, (5) is
Theorem~\ref{thm:dichot}: the scan finds no other non-triangle
non-smoothable face on any both-sides unit polytope.
\end{proof}

\begin{remark}[a threefold with no smoothing whose mirror has only
locally smoothable singularities]\label{rem:frozen}
Theorem~\ref{thm:pairbody} pairs a Calabi--Yau threefold that
cannot be smoothed (its $\FF_1$-cone point is rigid, with reduced
miniversal base a point,
the local model of Gross's exceptional case) with a mirror all of whose
singular points move locally.  Whether these local deformations assemble
to a global smoothing of $X^\circ$ is
exactly the frontier of \cite{paper2}: all $26+2+2$ of its singular points
sit over faces with $\ell(F^\circ) = 1$, so the single-face smoothing
techniques of \cite{paper2} do not apply.  The node subsystem fails the
all-nonzero relation of the nodal problem (Remark~\ref{rem:node-subsystem}),
while the four del Pezzo-cone points are the $\ell = 1$ germs that
\cite[Prop.~6.1, Rem.~6.2]{paper2} place beyond Gross's theorems, with no
Friedman-type obstruction known.  The mixed local-to-global problem remains
open.  The pair thus concentrates,
on two mirror-dual threefolds, every open smoothability phenomenon of this
series.
\end{remark}

\begin{remark}[the node subsystem]\label{rem:node-subsystem}
Let $q_1,\dots,q_{26}$ be the vertices of
$\Delta_{\FF_1}^\circ$.  For each of its $26$ unit-square faces
$S=(q_a,q_b,q_c,q_d)$ in cyclic order, form the diagonal-relation vector
\[
 \rho_S=e_a-e_b+e_c-e_d\in
 \mathbb{Q}^{\operatorname{Vert}(\Delta_{\FF_1}^\circ)}.
\]
Exact row reduction of the resulting $26\times26$ matrix gives rank $18$.
Exactly two rows are coloops.  The corresponding square faces, displayed
in cyclic order, are
\[
\begin{aligned}
S_-={}&\conv\{(-1,-1,0,-1),\ (-1,0,0,-1),\
                 (0,1,0,-1),\ (0,0,0,-1)\},\\
S_+={}&\conv\{(-1,-1,0,1),\ (0,-1,0,1),\
                 (0,0,0,1),\ (-1,0,0,1)\}.
\end{aligned}
\]
Deleting either row lowers the rank to $17$, deleting both lowers it to
$16$, and deleting any other row leaves rank $18$.
Thus every linear dependence among the $26$ node rows has coefficient zero
on both $\rho_{S_-}$ and $\rho_{S_+}$; in particular there is no dependence
with all coefficients non-zero.  The calculation is exact and is reproduced
by the ancillary row-reduction script.
\footnote{\texttt{paper3\_node\_relations.py}.}

For a hypothetical all-nodal threefold this fails the
Friedman--Batyrev--Kreuzer condition recalled in \cite[\S7]{paper2}.
It does \emph{not} decide the actual $X^\circ$: its four additional
del Pezzo-cone points have divisorial crepant resolutions, and their local
smoothing directions may contribute to the mixed local-to-global
obstruction map.  What the computation proves is that a purely nodal
cross-face rescue is impossible; any smoothing of $X^\circ$ must couple
the node directions to the del Pezzo directions.
\end{remark}

\section{Questions}\label{sec:questions}

\begin{question}\label{q:why}
Why triangles, zonotopes, and reflexive polygons?  Theorem~\ref{thm:class}
is an exhaustive fact about dimension four; a conceptual, dimension-free
proof (a local characterization of the unit-edge polygons admitting an
all-unit, all-transversally-smooth reflexive completion) would explain
it.  The failed local balancing statement of \S\ref{sec:refuted} shows the
answer is subtler than the star of the face; the appearance of exactly the
\emph{reflexive} polygons as the non-zonotope exceptions is suggestive in a
polar-duality-symmetric problem.
\end{question}

\begin{question}\label{q:mirror-smooth}
Is the mirror $X^\circ$ of Theorem~\ref{thm:pairbody} smoothable?  All its
germs are of type \textup{(S)} with $\ell(F^\circ)=1$.  A positive answer
would exhibit a Batyrev mirror pair with one member rigid against smoothing
and the other smoothable --- a maximally asymmetric pair for the
deformation--resolution exchange of mirror symmetry.  The node subsystem
alone admits no all-nonzero relation (Remark~\ref{rem:node-subsystem}); can
the four del Pezzo deformation directions supply the required mixed
local-to-global cancellation?
\end{question}

\begin{question}\label{q:mutable}
Corti--Filip--Petracci conjecturally relate smoothing components of $C(Q)$ to $0$-mutable
Laurent polynomials with Newton polygon $Q$ \cite{CFP}, and Filip
\cite{Filip25} constructs distinguished formal deformations in the
non-isolated case, detecting when their general fibres are smooth.  In that language, what
distinguishes the twelve face classes of Table~\ref{tab:occ} --- and what is
the mirror-theoretic meaning of the unique $\FF_1$ face?  One would like to
read Theorem~\ref{thm:dichot} off the canonical scattering diagram of the
mirror.
\end{question}

\begin{thebibliography}{Wue26b}

\bibitem[Alt97]{Altmann97}
K.~Altmann,
\emph{The versal deformation of an isolated toric Gorenstein singularity},
Invent.\ Math.\ \textbf{128} (1997), 443--479.

\bibitem[Bat94]{Batyrev94}
V.~V. Batyrev,
\emph{Dual polyhedra and mirror symmetry for Calabi--Yau hypersurfaces in
toric varieties},
J.\ Algebraic Geom.\ \textbf{3} (1994), 493--535.

\bibitem[Bra12]{Braun12}
V.~Braun,
\emph{The 24-cell and Calabi--Yau threefolds with Hodge numbers $(1,1)$},
J.\ High Energy Phys.\ \textbf{05} (2012) 101; arXiv:1102.4880.

\bibitem[CFP22]{CFP}
A.~Corti, M.~Filip, and A.~Petracci,
\emph{Mirror symmetry and smoothing Gorenstein toric affine 3-folds},
Facets of Algebraic Geometry vol.~I, LMS Lecture Note Ser.\ \textbf{472},
CUP, 2022, 132--163.

\bibitem[Fil26]{Filip25}
M.~Filip,
\emph{Laurent polynomials and deformations of non-isolated Gorenstein toric
singularities},
Adv.\ Math.\ \textbf{487} (2026), 110775; arXiv:2504.04486.

\bibitem[Gro97]{Gross97}
M.~Gross,
\emph{Deforming Calabi--Yau threefolds},
Math.\ Ann.\ \textbf{308} (1997), 187--220; updated version
arXiv:alg-geom/9506022v2 (2025).

\bibitem[KS00]{KreuzerSkarke00}
M.~Kreuzer and H.~Skarke,
\emph{Complete classification of reflexive polyhedra in four dimensions},
Adv.\ Theor.\ Math.\ Phys.\ \textbf{4} (2000), 1209--1230.

\bibitem[Wue26a]{paper1}
B.~J. Wuebben,
\emph{Non-smoothable Calabi--Yau threefolds from reflexive polytopes},
preprint, 2026; available at
\url{https://github.com/bwuebben/calabi-yau-smoothability}.

\bibitem[Wue26b]{paper2}
B.~J. Wuebben,
\emph{Smoothing Calabi--Yau threefolds in Gorenstein toric Fano fourfolds},
preprint, 2026; available at
\url{https://github.com/bwuebben/calabi-yau-smoothability}.

\end{thebibliography}

\bigskip
\noindent{\sc Bernd Johannes Wuebben, New York, NY,} \texttt{wuebben@gmail.com}

\end{document}
